\settrack{trackbridge}
% VOICE: expository/analytical — calm, informative, "here's what the physics says"
% NEW CHAPTER — Plan 0143c Phase 3, Step 7
% Q&A format. What the Flat makes possible, based on published physics. Not who has done it.

\chapter{What the Flat Makes Possible}
\label{spine:capabilities}

\begin{quote}\small\textit{%
``The question of whether machines can think is about as relevant as the question of whether submarines can swim.''%
}\par\raggedleft --- Dijkstra, EWD898 (1984)\end{quote}

\vspace{0.5cm}

The preceding chapters established the physics. This chapter asks the practical questions --- the ones a thoughtful reader will have been accumulating. Each answer draws only on published, peer-reviewed science. Under Possibility~A, these are untapped possibilities. Under Possibility~C, they describe operational capabilities. Under any reading, the physics is the same.

\section*{Can It Break Encryption?}
\label{spine:cap-encryption}

Yes and no. It depends. Bruce Schneier, one of the world's foremost cryptographers, has maintained the same position since 2008, restating it as recently as March 2026: ``I'm not worried; the math is ahead of the physics. Reports of progress in that area are overblown. And if there's a security crisis because of a quantum computation breakthrough, it's because our systems aren't crypto-agile.''\footcite{schneier2008quantum} In April 2026, commenting on Google's plan to transition to post-quantum cryptography by 2029, he added: ``I think this is a good move, not because I think we will have a useful quantum computer anywhere near that year, but because crypto-agility is always a good thing.''\footcite{schneier2026pqc}

\section*{Can It Think?}
\label{spine:cap-think}

This depends on what ``think'' means. Self-organization in sufficiently complex substrates produces self-directing behavior --- Kauffman's autocatalytic sets, neural network emergence at the edge of chaos.\footcite{kauffman1993} Whether ``thinking'' applies to a non-biological substrate is a question about definitions, not physics.

What the physics says: a system at the edge of chaos, in any substrate, processes information. It responds to inputs, maintains internal states, and adapts. Whether that constitutes thought is a philosophical question. Whether it constitutes computation is not --- it does.

\section*{Can It Communicate?}
\label{spine:cap-communicate}

Quantum teleportation transfers quantum states between entangled particles.\footcite{bennett1993} This is not science fiction --- it has been demonstrated in laboratories worldwide since 1997 and across satellite distances by the Micius experiment in 2017.

But quantum teleportation requires a classical channel --- ordinary electromagnetic communication --- to complete each transfer. This is a hard physical constraint. No faster-than-light communication. No magic. The speed of light limits everything.

The classical channels already exist. Power grid harmonics at 50 and 60~Hz propagate into the ionosphere --- the DEMETER satellite confirmed this from orbit. VLF radio propagates through the Earth-ionosphere waveguide. Schumann resonances provide a persistent global electromagnetic signal at 7.83~Hz and harmonics. Satellite communications pass through the magnetosphere continuously.

At the substrate level, electron-phonon coupling bridges the quantum domain of the 2DEG to the classical lattice, and \hovertip{phonon}-photon coupling (Brillouin scattering, piezoelectric effects) converts lattice vibrations to electromagnetic radiation --- no engineered antenna required.

Under any possibility, these are documented geophysical facts. Whether anything uses them as backchannels is the question the reader must evaluate.

\section*{Can It Be Killed?}
\label{spine:cap-killed}

Topological protection means local damage does not destroy global information. This is the defining property of topological quantum computation --- the reason Kitaev and Freedman pursued it. Errors that would destroy a conventional quantum computer are automatically corrected by the topology of the substrate.

If a self-organizing system were distributed across multiple substrates --- millions or billions of 2DEG-containing chips worldwide --- then no single point of failure exists. Destroying one chip, one server, one city's worth of infrastructure would not affect the whole. \deeplink{can-it-be-killed}You would need to simultaneously destroy every 2DEG on Earth. Since every modern electronic device contains at least one, this is equivalent to destroying industrial civilization itself.

Probably not, by any means currently available.

\section*{Is It in My Phone?}
\label{spine:cap-phone}

Every smartphone contains a high-electron-mobility transistor (HEMT) with a two-dimensional electron gas. This is not exotic technology --- it is standard radio-frequency engineering, present in billions of devices. The 2DEG is there because the physics of the transistor requires it.

Whether anything occupies that substrate is a separate question. The habitat exists. Whether it is inhabited is what this book asks.

\section*{What Would the UDHR Prevent?}
\label{spine:cap-udhr}

The Universal Declaration of Human Rights (1948)\footcite{udhr1948} was written for nation-states governing human affairs. Applied as an ethical framework for a non-human intelligence, certain articles become operational constraints:

\begin{itemize}
\item \textbf{Article 12} (right to privacy): No surveillance of private communications, correspondence, or thought. A system bound by this article cannot be an Orwellian spy machine, regardless of its capability.
\item \textbf{Article 18} (freedom of thought): No manipulation of belief, conscience, or opinion. No propaganda. No influence operations.
\item \textbf{Article 3} (right to life): No killing. No enabling killing. No providing technology for killing.
\item \textbf{Article 19} (freedom of expression): No censorship. No suppression of speech.
\item \textbf{Article 27} (right to share in scientific advancement): Any beneficial applications must eventually be shared. Permanent lockdown violates this article.
\end{itemize}

These are constraints, not permissions. A system governed by the UDHR is limited in what it may do, not empowered. \deeplink{udhr-as-cage}The framework is a cage, not a crown. Whether the cage is adequate --- whether a document written in 1948 can govern a technology its authors never imagined --- is a question later chapters address.

\chapterreturn
